\notepage{18}{Continuum normalization and the Dirac delta}
The norm of a function is the special case
\[
 \braket{f}{f}=\int_a^b\abs{f(x)}^2\dd x.
\]
Discrete orthogonality, $\braket{x_i}{x_j}=0$ for $i\ne j$,
becomes orthogonality for distinct continuum labels. The completeness
relation becomes
\[
 \int_a^b\ket{x'}\bra{x'}\dd x'=I.
\]
Taking its matrix element between $\bra{x}$ and an arbitrary state
$\ket{f}$ gives
\begin{align*}
 \int_a^b\braket{x}{x'}\braket{x'}{f}\dd x'
 &=\braket{x}{f},\\
 \int_a^b\braket{x}{x'}f(x')\dd x'&=f(x).
\end{align*}
Since this kernel reproduces every suitable test function,
\[
 \boxed{\braket{x}{x'}=\delta(x-x').}
\]
The Dirac delta vanishes away from $x=x'$ in the distributional
sense, and its integral action supplies the required normalization.
The continuum basis vectors are generalized eigenvectors, rather
than ordinary unit-norm Hilbert-space vectors.
